\chapter{Conclusiones}\label{cap:conclusiones}

Este trabajo probó cómo transferir lo que se sabe de una persona en la Tierra a su vida en el espacio, usando el reposo
en cama de la NASA como análogo y validando por sujetos. Las tres preguntas del \cref{cap:intro} tienen estas
respuestas.

\begin{description}
\item[P1. ¿Sirve el monitor de Tierra?] No en el peso, que es la variable que más se desplaza en cama: deja fuera el
  \cifra{peso.tierra.fa}\pct{} de los días normales. En las demás variables las falsas alarmas son menores, y su
  exceso sobre el piso que tiene la misma regla dentro del pre-reposo es chico y no se distingue bien del ruido.
\item[P2. ¿Cuánto mejora un rango adaptado?] Con las mismas falsas alarmas que el monitor de Tierra calibrado, el rango
  adaptado es más angosto en todas las variables. Detecta más escalones de forma clara en el peso,
  \cifra{peso.rec.dif} puntos \cifra{peso.rec.dif.ic}, y en la presión diastólica, \cifra{pa_dia.rec.dif}
  \cifra{pa_dia.rec.dif.ic}. En el pulso (\cifra{pulso.rec.dif} \cifra{pulso.rec.dif.ic}) y la presión sistólica
  (\cifra{pa_sis.rec.dif} \cifra{pa_sis.rec.dif.ic}) la mejora es menor y su intervalo cruza el cero. Detecta en los
  dos sentidos, a favor y en contra de la deriva del reposo. Lo aprendido de otras personas sirve sobre todo cuando
  aún no hay días propios y en las variables que más cambian; con 28 días propios, la mayor parte de la mejora viene de
  la propia persona.
\item[P3. ¿El pre-entrenamiento en Tierra mejora el pronóstico?] El PatchTST adaptado le gana por poco a la media de
  7 días (\cifra{par.todo.media7.r} ajustando todo). Frente al mismo modelo entrenado desde cero la diferencia es de
  \cifra{par.todo.cero.r} o menos, con una sola semilla y brazos entrenados en condiciones distintas. Todavía no se
  puede decir si el pre-entrenamiento aporta.
\end{description}

La auditoría confirmó que las cifras son las del kernel y que no hay fuga entre sujetos, y obligó a decir con más
precisión lo que el trabajo muestra. La lección más útil para NOVA es que, con pocos datos y un cambio que es sobre
todo de nivel, la transferencia que funciona es la más simple: la línea base de cada persona, la forma típica del
cambio aprendida de otras y los días propios que se van juntando, con un ancho calibrado y una regla fija y revisable.
Es también la que mejor calza con la regla del proyecto, porque son pocos números que se pueden revisar y reajustar a
bordo.

Queda por resolver si una red pre-entrenada en Tierra agrega algo encima de eso. La corrida v4 está diseñada para
responderlo con el pulso y la presión arterial, en igualdad de condiciones y con varias semillas.
